% Generated draft: requires mathematical and source review.
\documentclass[11pt]{article}
\usepackage[margin=25mm]{geometry}
\usepackage{fontspec}
\setmainfont{Latin Modern Roman}
\usepackage{amsmath,amssymb,mathtools,unicode-math}
\setmathfont{Latin Modern Math}
\usepackage{enumitem,needspace,xurl,hyperref}
\setlength{\parindent}{0pt}
\setlength{\parskip}{5pt}
\setlength{\emergencystretch}{3em}
\newcommand{\N}{\mathbb{N}}
\newcommand{\Z}{\mathbb{Z}}
\newcommand{\Q}{\mathbb{Q}}
\newcommand{\R}{\mathbb{R}}
\newcommand{\C}{\mathbb{C}}
\newcommand{\F}{\mathbb{F}}
\newcommand{\A}{\mathbb{A}}
\newcommand{\PP}{\mathbb{P}}
\newcommand{\E}{\mathbb{E}}
\newcommand{\eps}{\varepsilon}
\newcommand{\dd}{\,\mathrm{d}}
\begin{document}
\section*{UnsolvedMath ID: 2308023}
\textbf{Source problem code:} catalogue entry
\subsection{Definitions and mathematical statement}
Identify \(\C\) with \(\R^2\), write \(\dd A\) for area measure, and set \(\overline\partial=\tfrac12(\partial_x+i\partial_y)\). For a Borel set \(\Omega\subseteq\C\), its characteristic function defines a distribution, with
\[
\langle\overline\partial\chi_\Omega,\varphi\rangle
=-\int_\Omega\overline\partial\varphi\,\dd A
\qquad(\varphi\in C_c^\infty(\C)).
\]
For \(0<\alpha<1\), a customary inhomogeneous H\"older norm is
\[
\|f\|_{\operatorname{Lip}\alpha}
=\|f\|_\infty+
\sup_{z\ne w}\frac{|f(z)-f(w)|}{|z-w|^\alpha}.
\]
O'Farrell asks to characterize the Borel sets \(\Omega\) for which, at a fixed intended exponent \(\alpha>0\), a finite constant \(C_\Omega\) satisfies
\[
|\langle\overline\partial\chi_\Omega,f\rangle|
\le C_\Omega\|f\|_{\operatorname{Lip}\alpha}.
\]
At minimum this asks for the bound on smooth test functions; applying the pairing to all H\"older functions requires a specified extension convention. The source's full exponent range is \(\alpha>0\), not just \((0,1)\): for \(\alpha\ge1\), the intended higher-order H\"older/Zygmund convention must be stated rather than reusing a first-difference definition that would change the class.

The source reformulates the request as asking when \(\chi_\Omega\) acts on \(B^{\alpha-1}_{\infty,\infty}\), and singles out Swiss-cheese sets. Here \(B^s_{\infty,\infty}\) conventionally denotes a Besov space of distributions of smoothness \(s\), with norm given by dyadic pieces \(\sup_{j\ge0}2^{js}\|\Delta_j u\|_\infty\) in an inhomogeneous decomposition. The excerpt does not specify that convention or whether ``acts'' means a continuous linear functional or a multiplication operator; those are different properties. A common Swiss-cheese model is a closed disk with a countable union of open disks removed, but the allowed holes and geometry are unspecified here.

The motivation is to extend the Cauchy--Pompeiu formula beyond weakly differentiable functions: its boundary Cauchy integral can be formed for continuous data, while its \(\overline\partial f\) area term requires an appropriate interpretation. The formula's domain and boundary regularity and the precise equivalence with the Besov condition are not supplied, so no unqualified equivalence theorem is asserted. The supplied OPEN-TRIAGE summary reports open status in the 2018 edition and no definitive resolution located.
\subsection{Short English statement}
Which planar Borel sets have a characteristic-function derivative that acts boundedly on a H\"older class, permitting a broader Pompeiu formula? The Hölder exponent conventions, Besov interpretation, and Swiss-cheese geometry need to be fixed.
\subsection{Sources}
\begin{itemize}
\item Local UnsolvedMath registry, permanent ID \(2308023\), exact source code \texttt{catalogue entry}, ``Research Problems in Function Theory --- Problem 8.23.''
\item Hayman, \emph{Research Problems in Function Theory} (2018), Problem 8.23, attributed to A. G. O'Farrell; \url{https://arxiv.org/abs/1809.07200}. Source-TeX extraction; recorded access 2026-07-29.
\item The source mentions Pompeiu's formula and Besov reformulation without a separate precise reference. The supplied OPEN-TRIAGE summary provides no authoritative current-resolution citation.
\end{itemize}
\end{document}
